\documentclass[10pt,a4paper]{amsart}
\usepackage[margin=19mm]{geometry}
\usepackage{amsmath,amssymb,mathtools}
\usepackage[scr=rsfs,cal=euler]{mathalfa}
\usepackage{xfrac}
\usepackage[unicode,hidelinks,bookmarks=false]{hyperref}
\newcommand{\ie}{i.e.\ }
\newcommand{\cf}{cf.\ }
\newcommand{\resp}{respectively\ }
\newcommand{\Subsection}{\subsection}
\providecommand{\nobreakdash}{\nobreak\hbox{-}}
\setlength{\emergencystretch}{2em}
\multlinegap0pt
\usepackage{amssymb}
\usepackage{mathtools}
\usepackage[shortlabels]{enumitem}
\newtheorem{lemma}{Lemma}
\newtheorem{theorem}{Theorem}
\newtheorem{proposition}{Proposition}
\usepackage{cleveref}
\crefname{lemma}{Lemma}{Lemmas}
\crefname{theorem}{Theorem}{Theorems}
\crefname{proposition}{Proposition}{Propositions}
\newcommand{\DG}{\Delta}
\newcommand{\Ent}{H}
\newcommand{\KL}{D_{\mathrm{KL}}}
\newcommand{\R}{\mathbb R}
\newcommand{\cc}{\mathbf c}
\newcommand{\e}{\mathrm e}
\newcommand{\ee}{\mathbf e}
\newcommand{\kk}{\mathbf k}
\newcommand{\pot}{\psi}
\newcommand{\qq}{\mathbf q}
\newcommand{\ssv}{\mathbf s}
\newcommand{\uu}{\mathbf u}
\title{An Improved Upper Bound for Multicolour Ramsey Numbers}
\author{Sunghyeon Jo}
\date{Source: arXiv:2609.04596v1 (CC BY 4.0)}
\begin{document}
\maketitle

\noindent\emph{Source and licence.} An Improved Upper Bound for Multicolour Ramsey Numbers. Version arXiv:2609.04596v1 (CC BY 4.0), distributed by arXiv under the Creative Commons Attribution 4.0 International licence, http://creativecommons.org/licenses/by/4.0/, as recorded in the arXiv licence metadata for this exact version and shown on the abstract page https://arxiv.org/abs/2609.04596v1. Full licence notice: \texttt{licenses/arxiv-2609.04596-CC-BY-4.0.txt}.\par\smallskip

\noindent\textit{Editorial scope.} Counted unit: the article's proposition thm:offdiag-flex (source0.tex lines 574--588). The article's complete original proof of it is delivered with it (source0.tex lines 754--901, one \texttt{proof} environment of 148 lines, which the article places away from the statement). No mathematics is written, repaired or re-derived: the statement and the proof are the article's own lines, in the article's own notation, and no retained line is substituted or re-flowed.  Retained uncounted as the source-local context the proof needs: lines 186--191; lines 198--216; lines 206--210; lines 212--215; lines 241--250; lines 278--292; lines 283--285; lines 288--291; lines 328--331; lines 360--398; lines 380--386; lines 549--552; lines 591--595; lines 597--633; lines 604--608; lines 624--626; lines 637--640; lines 701--716; lines 709--715.  Nothing was omitted from any retained span, so the package carries no omission record.  No retained line contains a citation, so the wrapper carries no bibliography and no bibliography entry.  No retained line contains a figure environment, an image-inclusion command or drawing code, so no image is delivered and the package declares no asset dependency.  Editorial formatting only, in addition: an AMS \texttt{amsart} preamble that restates the article's own declarations and macros used by the retained lines, namely the environments \texttt{lemma}, \texttt{theorem}, \texttt{proposition} on separate counters, together with the standard packages those lines need.  The retained environments print in this document as Lemma 1, Theorem 1, Proposition 1 in order of appearance, whereas the article prints the same items with the numbering of its own full document.  The complete native source of the article as distributed in the arXiv e-print for this exact version is bundled unchanged as \texttt{sources/209375/source0.tex}.
\begin{equation}\label{eq:ES}
 R(k_1,\ldots,k_r)
 \le \binom{B-r}{k_1-1,\ldots,k_r-1}
 \le \binom{B}{k_1,\ldots,k_r}
 \le \exp(\Ent(\kk)).
\end{equation}

\begin{lemma}\label{lem:entropy}
Let $\kk\in\R_{>0}^r$ and let $\uu\in\R_{\ge0}^r$ satisfy
$\cc=\kk-\uu\in\R_{>0}^r$.  Put
\[
 B=\sum_i k_i,\qquad U=\sum_i u_i,
 \qquad q_i=\frac{k_i}{B}.
\]
Then
\begin{equation}\label{eq:entropy-identity}
 \Ent(\kk)+\sum_i u_i\log q_i-\Ent(\cc)
 =(B-U)\,\KL\!\left(\frac{\cc}{B-U}\middle\|\qq\right)
 =:\DG_{\kk}(\uu).
\end{equation}
Moreover, if $z_i=u_i-q_iU$, then
\begin{equation}\label{eq:entropy-lower}
 \DG_{\kk}(\uu)\ge
 \frac12\sum_{i=1}^r\frac{z_i^2}{k_i}.
\end{equation}
\end{lemma}

\begin{equation}
 \Ent(\kk)+\sum_i u_i\log q_i-\Ent(\cc)
 =(B-U)\,\KL\!\left(\frac{\cc}{B-U}\middle\|\qq\right)
 =:\DG_{\kk}(\uu).
\end{equation}

\begin{equation}
 \DG_{\kk}(\uu)\ge
 \frac12\sum_{i=1}^r\frac{z_i^2}{k_i}.
\end{equation}

\begin{lemma}\label{lem:potential}
Let $\kk\in\R_{>0}^r$ satisfy $\pot(\kk)>0$, let
$\uu\in\R_{\ge0}^r$ and $\cc=\kk-\uu\in\R_{>0}^r$, and put
$U=\sum_i u_i$ and $m=\min_i k_i$.  For every $\Gamma>0$,
\[
 \Gamma\bigl(\pot(\kk)-\pot(\cc)\bigr)
 \le \frac{4\Gamma U}{r}
      +\frac12\DG_{\kk}(\uu)+10\Gamma^2m.
\]
\end{lemma}

\begin{lemma}\label{lem:weighted}
Let $\qq=(q_1,\ldots,q_r)$ be a probability vector with every $q_i>0$, let
$0<\eta<\min_iq_i$, and let $s_\ast$ be a positive integer.  Every
$r$-edge-coloured $K_n$ contains pairwise disjoint sets
$S_1,\ldots,S_r,W$ such that, writing $s_i=|S_i|$ and $s=\sum_i s_i$,
\begin{equation}\label{eq:weighted-size}
 |W|\ge n(1+\eta)^s\prod_{i=1}^r q_i^{s_i},
\end{equation}
$S_i$ is a colour-$i$ clique, and every edge between $S_i$ and $W$ has colour
$i$.  Moreover, either $s=s_\ast$, or
\begin{equation}\label{eq:weighted-degree}
 |N_i(w)\cap W|\ge(q_i-\eta)|W|-1
 \qquad(w\in W,\ i\in[r]).
\end{equation}
\end{lemma}

\begin{equation}
 |W|\ge n(1+\eta)^s\prod_{i=1}^r q_i^{s_i},
\end{equation}

\begin{equation}
 |N_i(w)\cap W|\ge(q_i-\eta)|W|-1
 \qquad(w\in W,\ i\in[r]).
\end{equation}

\begin{equation}\label{eq:corr-params}
 \beta_r=\frac1{4r2^{r-1}(r+1)},
 \qquad C_{r,d}\le C_0rd^2\log(2rd),
\end{equation}

\begin{theorem}\label{thm:book}
Let $r\ge2$ and $d\ge3$ be integers, and let $\beta_r$ and $C_{r,d}$ be as in
\eqref{eq:corr-params}.  Let $t,m_1,\ldots,m_r$ be positive integers and let
$\lambda_0>0$.  For each
$i\in[r]$, let $p_i,\delta_i>0$ and define
\[
 L_i=\log(1/\delta_i),\qquad L=\max_iL_i,
 \qquad L_{p,i}=\log(1/p_i),
\]
\[
 \Pi_i=\frac{3\delta_i}{p_i}
       +\frac{6L_iL_{p,i}}{\lambda_0},
\]
and
\[
 \rho=\log(r/\beta_r),\qquad
 \Xi=2\rho+4C_{r,d}\lambda_0^{1/d}
      +\frac{12C_{r,d}L}{\lambda_0^{(d-1)/d}}.
\]
Suppose
\begin{equation}\label{eq:book-parameters}
 0<\delta_i\le\min\{p_i/4,1/4\},
 \qquad
 \lambda_0\ge\max\{2,6L\},
 \qquad
 t\ge\lambda_0(\min_i\delta_i)^{-1/(d-1)}.
\end{equation}
Let $X,Y_1,\ldots,Y_r$ be nonempty vertex sets such that
\begin{align}
 |N_i(x)\cap Y_i|&\ge p_i|Y_i|
 &&(x\in X,\ i\in[r]),\notag\\
 |Y_i|&\ge p_i^{-t}\e^{\Pi_it}m_i
 &&(i\in[r]),\notag\\
 |X|&\ge2rt\e^{rt\Xi}.\label{eq:book-reservoir}
\end{align}
Then there is a colour $i$, a colour-$i$ clique $T\subseteq X$ of size $t$,
and a set $P\subseteq Y_i\setminus T$ of size $m_i$ such that every edge
between $T$ and $P$ has colour $i$.
\end{theorem}

\begin{equation}
 0<\delta_i\le\min\{p_i/4,1/4\},
 \qquad
 \lambda_0\ge\max\{2,6L\},
 \qquad
 t\ge\lambda_0(\min_i\delta_i)^{-1/(d-1)}.
\end{equation}

\begin{equation}\label{eq:constant-choice}
 130\gamma+20\gamma^2+3\zeta+\frac6A<\frac1{576},
 \qquad \gamma<10^{-3}.
\end{equation}

\begin{equation}\label{eq:Gamma-eta}
 \Gamma=\gamma\theta_{r,d},
 \qquad
 \eta=\frac{16\Gamma}{r}.
\end{equation}

\begin{lemma}\label{lem:parameters}
There is an absolute constant $K_0$ such that the following holds.  Suppose
$\pot(\kk)>0$ and
\[
 m=\min_i k_i\ge m_0:=K_0\mathcal T_{r,d}.
\]
Let $q_i=k_i/B$ and define
\begin{equation}\label{eq:t-sstar}
 t=\lfloor\theta_{r,d}m\rfloor,
 \qquad
 s_\ast=\lceil4\Gamma rm\rceil,
\end{equation}
\[
 p_i=q_i-2\eta,
 \qquad
 \delta_i=\zeta p_i\theta_{r,d},
 \qquad
 L_i=\log(1/\delta_i),\quad L=\max_iL_i,
 \qquad
 \lambda_0=\frac{AL^2}{\theta_{r,d}}.
\]
Then the hypotheses in \eqref{eq:book-parameters} hold, and
\begin{equation}\label{eq:p-window}
 \frac1{3r}\le p_i\le\frac2r,
 \qquad
 L_i=\Theta(\log(2rd))\quad\text{uniformly in $i$},
\end{equation}
\begin{equation}\label{eq:Pi-bound}
 \Pi_i\le\left(3\zeta+\frac6A\right)\theta_{r,d}.
\end{equation}
Moreover, if $\Xi$ is as in \cref{thm:book}, then
\begin{equation}\label{eq:Xi-bound}
 rt\Xi+\log(2rt)\le\tfrac1{16}\,rm\log(2r).
\end{equation}
Finally, $\theta_{r,d}^2m\ge C_{\mathrm{rnd}}$ once $K_0$ is sufficiently
large.
\end{lemma}

\begin{equation}
 t=\lfloor\theta_{r,d}m\rfloor,
 \qquad
 s_\ast=\lceil4\Gamma rm\rceil,
\end{equation}

\begin{equation}
 \Pi_i\le\left(3\zeta+\frac6A\right)\theta_{r,d}.
\end{equation}

\begin{equation}\label{eq:q-window}
 \frac1{2r}<q_i<\frac2{r+1}
 \qquad(i\in[r]).
\end{equation}

\begin{lemma}\label{lem:reservoir}
Under the hypotheses and notation of \cref{lem:parameters}, let $N$ satisfy
\[
 \log N\ge \Ent(\kk)-\Gamma\pot(\kk),
\]
and apply \cref{lem:weighted} with weights $q_i=k_i/B$ and stopping size
$s_\ast$.  If the procedure stops with $s<s_\ast$, then, provided $K_0$ is
sufficiently large,
\begin{equation}\label{eq:reservoir-conclusion}
 \eta|W|\ge1,
 \qquad
 |W|\ge2rt\e^{rt\Xi},
 \qquad
 |W|\ge\max_i p_i^{-t}\e^{\Pi_it}.
\end{equation}
\end{lemma}

\begin{equation}
 \eta|W|\ge1,
 \qquad
 |W|\ge2rt\e^{rt\Xi},
 \qquad
 |W|\ge\max_i p_i^{-t}\e^{\Pi_it}.
\end{equation}

\begin{proposition}\label{thm:offdiag-flex}
There are absolute constants $c,C>0$ such that the following holds.  Let
$r\ge2$ and let $d$ be an integer with
\[
 3\le d\le\max\{3,\log(2r)\}.
\]
Then there is a constant $A_{r,d}\le C\theta_{r,d}\mathcal T_{r,d}$ such that,
for every positive integer vector $\kk=(k_1,\ldots,k_r)$,
\[
 R(\kk)\le
 \left\lceil
 \exp\bigl(\Ent(\kk)-c\theta_{r,d}\pot(\kk)+A_{r,d}\bigr)
 \right\rceil.
\]
\end{proposition}

\begin{proof}[Proof of \cref{thm:offdiag-flex}]
Fix $r,d$ and the constants above, and write
$\theta=\theta_{r,d}$.  Let
\[
 m_0=K_0\mathcal T_{r,d},
 \qquad
 A_{r,d}=\Gamma m_0.
\]
Since $\Gamma=\gamma\theta$, this satisfies the bound
$A_{r,d}\le C\theta_{r,d}\mathcal T_{r,d}$ of \cref{thm:offdiag-flex} with
$C=\gamma K_0$, and the saving coefficient there is $c=\gamma$.  We prove by
induction on $B=\sum_i k_i$ that
\begin{equation}\label{eq:induction}
 R(\kk)\le
 \left\lceil\exp\bigl(\Ent(\kk)-\Gamma\pot(\kk)+A_{r,d}\bigr)\right\rceil.
\end{equation}
If $\pot(\kk)=0$, this follows from \eqref{eq:ES}.  If
$m=\min_i k_i\le m_0$, then $\pot(\kk)\le m_0$ and the additive term again
makes \eqref{eq:induction} weaker than \eqref{eq:ES}.  We may therefore assume
\[
 \pot(\kk)>0,
 \qquad m>m_0.
\]
We record once how integer roundings are absorbed.  For every positive
integer vector $\mathbf x$ we have $\pot(\mathbf x)\le\min_ix_i$, while
writing $\Ent(\mathbf x)=\sum_ix_i\log\bigl(\sum_jx_j/x_i\bigr)$ shows that
the summand corresponding to a minimal coordinate is already at least
$(\min_ix_i)\log r$; since $\Gamma\le\gamma b<\log2$, every exponent
appearing in \eqref{eq:induction} is therefore nonnegative, and each ceiling
is at most twice the corresponding exponential.  Consequently, a margin of
$\log2$ in an exponent inequality below absorbs all integer roundings.  As
$\gamma$, $\zeta$ and $A$ are already fixed, we may choose
$C_{\mathrm{rnd}}$, and then $K_0$, so large that the margins obtained in
\eqref{eq:long-margin} and \eqref{eq:page-positive} below are at least
$\log2$; we use this without further comment.
Let $N$ be the right-hand side of \eqref{eq:induction}, and consider an
arbitrary $r$-colouring of $K_N$.  Apply \cref{lem:weighted} with
$q_i=k_i/B$, the parameter $\eta$ from \eqref{eq:Gamma-eta}, and the stopping
size $s_\ast$ from \eqref{eq:t-sstar}.  Let $S_i,W$ be the resulting sets and
write $s_i=|S_i|$, $\ssv=(s_1,\ldots,s_r)$ and $s=\sum_i s_i$.  If $s_i\ge k_i$ for some $i$ then we are
done, so assume $s_i<k_i$ for every $i$.

Suppose first that $s=s_\ast$.  Put $\cc=\kk-\ssv$ and
$\DG=\DG_{\kk}(\ssv)$.  By \eqref{eq:weighted-size} and
\cref{lem:entropy},
\begin{align*}
 \log|W|
 &\ge \Ent(\kk)-\Gamma\pot(\kk)+A_{r,d}
      +\sum_i s_i\log q_i+s\log(1+\eta)\\
 &=\Ent(\cc)-\Gamma\pot(\cc)+A_{r,d}
   +\DG+s\log(1+\eta)
   -\Gamma\bigl(\pot(\kk)-\pot(\cc)\bigr).
\end{align*}
Using \cref{lem:potential},
$\log(1+\eta)\ge\eta/2=8\Gamma/r$, and
$s=s_\ast\ge4\Gamma rm$, we obtain
\begin{align}
 &\DG+s\log(1+\eta)
 -\Gamma\bigl(\pot(\kk)-\pot(\cc)\bigr)\notag\\
 &\hspace{2cm}\ge
 \frac12\DG+\frac{4\Gamma s}{r}-10\Gamma^2m
 \ge6\Gamma^2m.\label{eq:long-margin}
\end{align}
Since $6\Gamma^2m\ge6\gamma^2C_{\mathrm{rnd}}\ge\log2$ by
\cref{lem:parameters} and the choice of $C_{\mathrm{rnd}}$, the rounding
convention gives $|W|\ge R(\cc)$ by the induction hypothesis.  Hence $W$
contains, for some $j$, a colour-$j$ clique of size $k_j-s_j$, which
together with $S_j$ completes the target $K_{k_j}$.

We may therefore assume $s<s_\ast$.  The minimum-degree conclusion
\eqref{eq:weighted-degree} holds, and \cref{lem:reservoir} gives
\eqref{eq:reservoir-conclusion}.  In particular, for every $w\in W$,
\[
 |N_i(w)\cap W|
 \ge(q_i-\eta)|W|-1
 \ge(q_i-2\eta)|W|=p_i|W|.
\]
For each colour $i$ with $s_i+t<k_i$, let
\[
 \cc^{(i)}=\kk-\ssv-t\ee_i,
 \qquad m_i=R(\cc^{(i)}).
\]
If $s_i+t\ge k_i$, set $m_i=1$; in that case a colour-$i$ spine of size
$t$ already completes the target.
Fix a colour $i$ with $s_i+t<k_i$.  Put
$\uu=\ssv+t\ee_i$, $U=s+t$, and
$\DG_i=\DG_{\kk}(\uu)$.  By the induction hypothesis,
\eqref{eq:entropy-identity} and the rounding convention, the page condition
$|W|\ge p_i^{-t}\e^{\Pi_it}m_i$ follows once
\begin{align*}
 \mathcal M_i:={}&
 \DG_i+s\log(1+\eta)
 -\Gamma\bigl(\pot(\kk)-\pot(\cc^{(i)})\bigr)\\
 &\qquad -t\log(q_i/p_i)-\Pi_it\ge\log2.
\end{align*}
In the notation of \cref{lem:entropy}, $z_i=s_i+t-q_i(s+t)$.  From
\eqref{eq:q-window},
$s<s_\ast\le5\Gamma rm$ (note that $\Gamma rm\ge\gamma C_{\mathrm{rnd}}\ge1$),
$t\ge\theta m/2$, and $\gamma<10^{-3}$, we have
\[
 z_i\ge(1-q_i)t-q_is\ge\frac t6.
\]
Since $k_i<2m$, \eqref{eq:entropy-lower} gives
\[
 \DG_i\ge\frac{z_i^2}{2k_i}\ge\frac{t^2}{144m}.
\]
Applying \cref{lem:potential} to the potential term of $\mathcal M_i$, and
using $s\log(1+\eta)\ge4\Gamma s/r$, gives
\[
 \mathcal M_i\ge
 \frac12\DG_i-\frac{4\Gamma t}{r}-10\Gamma^2m
 -t\log(q_i/p_i)-\Pi_it.
\]
Because $t\ge\theta m/2$,
\[
 \frac12\DG_i\ge\frac{\theta t}{576}.
\]
Moreover, $q_i\ge1/(2r)$ and $p_i=q_i-2\eta$ imply
\[
 \log(q_i/p_i)
 =-\log\left(1-\frac{2\eta}{q_i}\right)
 \le128\gamma\theta,
\]
and $\Pi_it\le(3\zeta+6/A)\theta t$ by \eqref{eq:Pi-bound}.  Thus
\begin{equation}\label{eq:page-positive}
 \mathcal M_i\ge
 \theta t\left(
 \frac1{576}-130\gamma-20\gamma^2-3\zeta-\frac6A
 \right)\ge\log2
\end{equation}
by \eqref{eq:constant-choice}, since
$\theta t\ge\theta^2m/2\ge C_{\mathrm{rnd}}/2$ and $C_{\mathrm{rnd}}$ is
sufficiently large.  Hence
\[
 |W|\ge p_i^{-t}\e^{\Pi_it}m_i
\]
for every colour $i$ with $s_i+t<k_i$.  If $s_i+t\ge k_i$, the same
inequality, now with $m_i=1$, is the last conclusion of \cref{lem:reservoir}.
We may now apply \cref{thm:book} with
$X=Y_1=\cdots=Y_r=W$.  It gives a colour $i$, a colour-$i$ spine $T$ of size
$t$, and a page set $P$ of size $m_i$.  If $s_i+t\ge k_i$, then
$S_i\cup T$ already contains a colour-$i$ $K_{k_i}$.  Otherwise
$|P|=R(\cc^{(i)})$, so $P$ contains, for some $j$, a colour-$j$ clique
whose size is the $j$th coordinate of $\cc^{(i)}$.  If $j=i$, join it to
$S_i\cup T$; if $j\ne i$,
join it to $S_j$.  In either case a target $K_{k_j}$ is completed.  This
proves \eqref{eq:induction}, and hence the proposition.
\end{proof}

\end{document}
